\begin{beginnerstatement}[Kurz erklärt: Programme im Entwicklungsbetrieb]

Ein Entwicklungsserver macht eine noch nicht ausgelieferte Anwendung während
der Arbeit lokal erreichbar. \texttt{run} startet dafür Vite und bei aktivem
Backend zusätzlich Uvicorn. Vite bedient das Frontend, Uvicorn das
FastAPI-Backend. Ein laufendes Programm heißt Prozess; ein Port ist dessen
nummerierter Zugangspunkt für Verbindungen. Im Vordergrund bleibt die Ausgabe
direkt im Terminal sichtbar und ein Abbruch beendet die erfassten Prozesse. Im
Hintergrund laufen sie abgekoppelt weiter. Eine PID identifiziert dabei einen
Prozess, während ein Log seine Meldungen festhält. \texttt{stop} beendet die
erfassten Prozesse und je nach Option auch weitere Listener auf ihren
konfigurierten Ports.

\end{beginnerstatement}
